\chapter{Algebraic Structures}

\emph{Algebraic structures} are mathematical systems consisting of a set equipped with one or more operations 
that satisfy certain axioms. 

\section{Internal And External Operations}

An \emph{internal composition law} is a binary operation that takes two elements from a set and returns another element in the same set. 
Formally, for a set \(S\) and operation \(\circ\), we have:

\[
  \circ: S \times S \to S
\]

An \emph{external composition law} involves a second set acting on the structure, such as scalar multiplication in vector spaces:

\[
  \circ: \field \times V \to V
\]

where \(\field\) is a field and \(V\) is a vector space.

\subsection{Properties of Operations}

Let \(\circ\) and \(\circledast\) be binary operations on a set \(S\). The most important properties include:

\begin{itemize}
    
    \item \emph{Associativity:} \((a \circ b) \circ c  = a \circ (b \circ c)\) for all \(a,b,c \in S\).
    
    \item \emph{Commutativity:} \(a \circ b = b \circ a\) for all \(a,b \in S\).
    
    \item \emph{Identity Element:} There exists \(e \in S\) such that \(a \circ e = e \circ a = a\) for all \(a \in S\).
    
    \item \emph{Inverse Element:} For every \(a \in S\), there exists \(a^{-1} \in S\) such that \(a \circ a^{-1} = a^{-1} \circ a = e\).
    
    \item \emph{Distributivity:} \(a \odot (b \oplus c) = (a \odot b) \oplus (a \odot c)\) and/or \((b \oplus c) \odot a = (b \odot a) \oplus 
    (c \odot a)\).

\end{itemize}

\section{Homomorphisms And Isomorphisms}

Let \((G, \oplus)\) and \((H, \oplus')\) be two algebraic structures.

A \emph{homomorphism} is a function \(\varphi: G \to H\) such that:
    
\[
  \varphi(a \oplus b) = \varphi(a) \oplus \varphi(b), \quad \forall a,b \in G
\]
    
An \emph{isomorphism} is a bijective homomorphism. If such a map exists, we say the structures are \emph{isomorphic}, 
written as \(G \cong H\).

\section{Lattice}

A \emph{lattice} \((L, \land, \lor)\) is an algebraic structure with two binary, associative, and commutative operations, which satisfy 
the following properties:

\[
    a \lor (a \land b) = a
\]

and 

\[
    a \land (a \lor b) = a
\]

for all \(a,b \in L\).
 
\section{Semigroup \texorpdfstring{\((S, \oplus)\)}{}}

A \emph{semigroup} is a set \(S\) equipped with a binary operation \(\oplus\) that is \emph{associative}. This means:

\[
    (a \oplus b) \oplus c = a \oplus (b \oplus c) \quad \text{for all } a, b, c \in S.
\]

\section{Monoid \texorpdfstring{\((M, \oplus)\)}{}}

A \emph{monoid} builds on a semigroup by adding an \emph{identity element} \(e\) such that:

\[
    a \oplus e = e \oplus a = a \quad \text{for all } a \in M.
\]

This structure is useful when an operation must have a ``do nothing'' element, like \(0\) for addition or \(1\) for multiplication.

\section{Group \texorpdfstring{\((G, \oplus)\)}{}}

A \emph{group} is a monoid where every element has an \emph{inverse}:

\[
    \text{For each } a \in G, \text{ there exists } a^{-1} \in G \text{ such that } a \oplus a^{-1} = a^{-1} \oplus a = e.
\]

\section{Abelian Group \texorpdfstring{\((A, \oplus)\)}{}}

An \emph{Abelian group} (or commutative group) is a group where the operation is \emph{commutative}:

\[
    a \oplus b = b \oplus a \quad \text{for all } a, b \in A.
\]

\section{Ring \texorpdfstring{\((R, \oplus, \odot)\)}{}}

A \emph{ring} is a set with two operations:

\begin{itemize}

    \item \((R, \oplus)\) is an Abelian group.

    \item \((R, \odot)\) is a semigroup (associative multiplication).
    
    \item Multiplication distributes over addition:
  
    \[
        a \odot (b \oplus c) = a \odot b \oplus a \odot c.
    \]

  \end{itemize}

\section{Commutative Ring \texorpdfstring{\((R, \oplus, \odot)\)}{}}

A \emph{commutative ring} is a ring where multiplication is also \emph{commutative}:

\[
    a \odot b = b \odot a \quad \text{for all } a, b \in R.
\]

This is the kind of ring most often encountered in basic algebra, such as the integers \(\Z\).

\section{Field \texorpdfstring{\((\field, \oplus, \odot)\)}{}}

A \emph{field} is a commutative ring where every nonzero element has a \emph{multiplicative inverse}:

\[
    (\field \setminus \{0\}, \odot) \text{ is an Abelian group}.
\]

Fields like \(\Q\), \(\R\), and \(\C\) allow division (except by zero), enabling the full range of arithmetic operations.

\subsection{Field Proofs}

\subsubsection{Inverses are well-defined:}

We know \(a + -a = 0\). Now let us say \(-a' \ne -a\) and \(a -a' = 0\).

\begin{align*}
    a + -a  &= a - a' \quad \mid - a \\ 
        -a &= -a'
\end{align*}

This contradicts our initial assumption; therefore, they are the same. The proof is analogous for the multiplication:

\begin{align*}
    a \cdot a^{-1} &= a \cdot a'^{-1} \quad \mid * a^{-1} \\ 
    a^{-1}  &= a'^{-1} 
\end{align*}

\QED

\subsubsection{Neutral Element for the Addition is well-defined:}

We know \(a + 0 = a\). Now let us say \(0' \ne 0\) and \(a + 0' = a\).

\begin{align*}
  a + 0 = a + (a - a) &=  a + 0' \\ 
          a - a + a   &= a + 0' \\
           a   &= a + 0' \quad \mid -a \\ 
           0 &= 0'
\end{align*}

There is a contradiction to the assumption that there is more than one neutral element for the addition. Thus, 
it is well-defined.

Neutral Element for the Multiplication is well-defined: This proof is 
very similar to the previous one; thus, it will be skipped. 

\QED

\subsubsection{Minus times Minus is Plus:}

We will prove that \(ab = (-a)(-b)\) for any two elements of a field.

\begin{align*}
	ab &= (-a)(-b) \\
	ab + (-a)b &= (-a)(-b) + (-a)b \\ 
	b(a + (-a)) &= (-a)((-b) + b) = 0\\
\end{align*}

From that point, for the right side we can do the following:

\begin{align*}
	(-a)((-b) + b) &= 0\\
	(-a)(-b) + (-a)b &= 0\\
	(-a)(-b) &= -(-a)b\\
\end{align*}

Then, by looking at the left side of the original equality:

\begin{align*}
	b(a + (-a)) &= 0\\
	ab + (-a)b &= 0\\
	ab &= -(-a)b\\
\end{align*}

Thus, we have obtained our equality \(ab = (-a)(-b)\).

\QED

\section{Vector Space \texorpdfstring{\((V, \oplus, \cdot)\)}{}}

A \emph{vector space} over a field \(\field\) consists of:

\begin{itemize}

    \item An Abelian group \((V, \oplus)\) for vector addition.

    \item A scalar multiplication \(\cdot: \field \times V \to V\) satisfying:

    \begin{itemize}

        \item Distributivity over vector addition: \(a \cdot (v_1 + v_2) = a \cdot v_1 + a \cdot v_2\)
    
        \item Distributivity over field addition: \((a + b) \cdot v = a \cdot v + b \cdot v\)

        \item Associativity: \(a \cdot (b \cdot v) = (ab) \cdot v\)
    
        \item Identity: \(1 \cdot v = v\)

  \end{itemize}

\end{itemize}

Vector spaces provide the foundation for linear algebra, where scalars from a field combine with vectors from a set.

